\subsection{Linear Changes of Variables and Euclidean Invariance}

\begin{perspective} \label{per:coordinate-splitting-completed-product-measure} \hlt{Coordinate Splitting as a Product Identification}\\
For $n=m+k$, let $\Phi:\R^n\to\R^m\times\R^k$ be $\Phi(x_1,\ldots,x_n)=((x_1,\ldots,x_m),(x_{m+1},\ldots,x_n))$. Then $\Phi$ sends rectangles to product rectangles and preserves rectangular volume, so $m_n$ is the completion of $\Phi^{-1}_{\#}(m_m\otimes m_k)$. Equivalently, integration over $\R^{m+k}$ matches iterated integration over $\R^m\times\R^k$ whenever Fubini applies.\\
\hlt{Key Insight}: Coordinate splitting preserves Lebesgue mass.
\end{perspective}

\begin{proposition} \label{prop:lipschitz-images-lebesgue-measurable} \hlt{Lipschitz Images of Lebesgue Sets}\\
Let $\Psi:\R^n\to\R^n$ be Lipschitz and let $E\in\mathcal L^n$. Then $\Psi(E)\in\mathcal L^n$.\\
Hence every linear map $\R^n\to\R^n$ preserves Lebesgue measurability.
\begin{proof}
Let $L$ be a Lipschitz constant.
\begin{enumerate}[label=(\roman*)]
\setlength{\itemsep}{0pt}
\item $\forall I\subseteq\R^n$ bounded cells, $m_n^\ast(\Psi(I))\leq C_{n,L}\operatorname{vol}_n(I)$. Indeed, partition $I$ into grid cells of side at most $\delta$.\\
Each image has diameter at most $L\sqrt n\,\delta$, hence lies in a cube of volume at most $C_{n,L}\delta^n$. Summing and letting $\delta\downarrow0$ gives the estimate.
\item If $m_n^\ast(N)=0$ and $\varepsilon>0$, cover $N$ by cells $(I_k)$ with $\sum\limits_k\operatorname{vol}_n(I_k)<\varepsilon$.\\
Then $\Psi(N)\subseteq\bigcup\limits_k\Psi(I_k)$, so (i) gives $m_n^\ast(\Psi(N))\leq C_{n,L}\varepsilon$; hence $m_n^\ast(\Psi(N))=0$.
\item If $E$ is bounded, choose $F_\sigma$ set $F\subseteq E$ with $m_n^\ast(E\setminus F)=0$. Write $F=\bigcup\limits_jF_j$ with each $F_j$ compact.\\
Then $\Psi(F)=\bigcup\limits_j\Psi(F_j)$ is Borel, $\Psi(E\setminus F)$ is null by (ii), and $\Psi(E)=\Psi(F)\cup\Psi(E\setminus F)\in\mathcal L^n$.
\item For general $E$, use $E=\bigcup\limits_{r\geq1}(E\cap[-r,r]^n)$ and apply (iii) to each bounded term.
\end{enumerate}
\hlt{Structural Commentary}: Reduce to compact pieces plus null error.\\
\hlt{Key Insight}: Lipschitz maps preserve null sets.
\end{proof}
\end{proposition}

\begin{figure}[H]
  \centering
  \begin{tikzpicture}[
  scale=0.90,
  >=Stealth,
  every node/.style={font=\scriptsize},
  axis/.style={->, gray!70, line width=0.45pt},
  original/.style={fill=blue!8, draw=blue!60!black, line width=0.75pt, fill opacity=0.62},
  image/.style={fill=red!8, draw=red!70!black, line width=0.75pt, fill opacity=0.50},
  motion/.style={->, gray!70, line width=0.65pt}
]
  \begin{scope}
    \draw[axis] (-0.20,0) -- (2.35,0);
    \draw[axis] (0,-0.15) -- (0,1.55);

    \draw[image] (0.30,0.28) rectangle (1.95,1.06);
    \draw[original] (0.30,0.28) rectangle (1.12,1.06);
    \draw[motion] (1.14,0.67) -- (1.82,0.67);
  \end{scope}

  \begin{scope}[xshift=3.25cm]
    \draw[axis] (-0.20,0) -- (2.35,0);
    \draw[axis] (0,-0.15) -- (0,1.55);

    \draw[image] (0.34,0.34) rectangle (0.94,1.38);
    \draw[original] (0.34,0.34) rectangle (1.38,0.94);
    \draw[motion] (1.33,1.15) .. controls (1.66,0.88) and (1.66,0.50) .. (1.33,0.36);
  \end{scope}

  \begin{scope}[xshift=6.50cm]
    \draw[axis] (-0.20,0) -- (2.35,0);
    \draw[axis] (0,-0.15) -- (0,1.55);

    \draw[image] (0.78,0.28) -- (1.66,0.28) -- (1.18,1.18) -- (0.30,1.18) -- cycle;
    \draw[original] (0.30,0.28) rectangle (1.18,1.18);
    \draw[motion] (1.20,1.02) -- (1.50,1.02);
  \end{scope}
\end{tikzpicture}

  \caption{\textcolor{blue!60!black}{Original} and \textcolor{red!70!black}{image} cells under scale, swap, and shear; Fubini checks each coordinatewise.}
  \label{fig:elementary-linear-change-maps}
\end{figure}

\begin{theorem} \label{thm:linear-change-of-variables-nonnegative} \hlt{Linear Change of Variables for Nonnegative Functions}\\
Let $T\in\LL(\R^n)$ be invertible and let $f:\R^n\to[0,+\infty]$ be Lebesgue measurable.
\begin{enumerate}[label=(\roman*)]
\setlength{\itemsep}{0pt}
\item $f\circ T$ is Lebesgue measurable.
\item $\int f\circ T\,dm_n=\abs{\det T}^{-1}\int f\,dm_n$.
\end{enumerate}
\begin{proof}
\begin{enumerate}[label=(\roman*)]
\setlength{\itemsep}{0pt}
\item Since $(f\circ T)^{-1}((a,+\infty])=T^{-1}(f^{-1}((a,+\infty]))$ and $T^{-1}$ is Lipschitz, Proposition \ref{prop:lipschitz-images-lebesgue-measurable} gives measurability.
\item The $n=1$ identity follows from $m_1(\alpha E+\beta)=\abs{\alpha}m_1(E)$.\\
Let $\mathcal M_+$ denote the measurable maps $\R^n\to[0,+\infty]$, set $J_S(h)\defeq\int h(Sx)\,dm_n(x)$.\\
Set $\mathcal A_n\defeq\{S\in GL(n):J_S(h)=\abs{\det S}^{-1}\int h\,dm_n\ \forall h\in\mathcal M_+\}$.\\
If $S_1,S_2\in\mathcal A_n$, then $J_{S_1S_2}(h)=\abs{\det S_2}^{-1}J_{S_1}(h)=\abs{\det(S_1S_2)}^{-1}\int h\,dm_n$, so $S_1S_2\in\mathcal A_n$.\\
The generators $D_j(c):x_j\mapsto cx_j$, $P_{j,j+1}:x_j\leftrightarrow x_{j+1}$, and $H_{jk}(c):x_j\mapsto x_j+cx_k$ lie in $\mathcal A_n$ for $c\neq0$ by the $n=1$ identity, Fubini, and slice translation invariance, with determinant factors $\abs c$, $1$, and $1$.\\
Row reduction gives $T\in\mathcal A_n$; hence $J_T(f)=\abs{\det T}^{-1}\int f\,dm_n$.
\end{enumerate}
\hlt{Structural Commentary}: Verify elementary row operations, then compose.\\
\hlt{Key Insight}: Only scalings change mass.
\end{proof}
\end{theorem}

\begin{corollary} \label{cor:invertible-linear-images-lebesgue-sets} \hlt{Invertible Linear Images of Lebesgue Sets}\\
Let $T\in\LL(\R^n)$ be invertible and let $E\subseteq\R^n$.
\begin{enumerate}[label=(\roman*)]
\setlength{\itemsep}{0pt}
\item $E\in\mathcal L^n\Longleftrightarrow T(E)\in\mathcal L^n$.
\item If $E\in\mathcal L^n$, then $m_n(T(E))=\abs{\det T}\,m_n(E)$.
\end{enumerate}
\begin{proof}
\begin{enumerate}[label=(\roman*)]
\setlength{\itemsep}{0pt}
\item Since $T$ and $T^{-1}$ are Lipschitz, Proposition \ref{prop:lipschitz-images-lebesgue-measurable} gives $E\in\mathcal L^n\Longleftrightarrow T(E)\in\mathcal L^n$.
\item For $E\in\mathcal L^n$, Theorem \ref{thm:linear-change-of-variables-nonnegative} applied to $\mathbf{1}_{T(E)}$ gives $m_n(E)=\int\mathbf{1}_{T(E)}(Tx)\,dm_n(x)=\abs{\det T}^{-1}m_n(T(E))$.
\end{enumerate}
\hlt{Structural Commentary}: Measurability is Lipschitz; mass follows from indicators.\\
\hlt{Key Insight}: Indicators turn integral change into set change.
\end{proof}
\end{corollary}

\begin{corollary} \label{cor:linear-l1-change-of-variables} \hlt{Linear Change of Variables for $L^1$ Functions}\\
Let $T\in\LL(\R^n)$ be invertible, let $f\in\mathscr L^1(m_n)$, and set $c_T\defeq\abs{\det T}^{-1}$.
\begin{enumerate}[label=(\roman*)]
\setlength{\itemsep}{0pt}
\item $f\circ T\in\mathscr L^1(m_n)$.
\item $\int f\circ T\,dm_n=c_T\int f\,dm_n$.
\end{enumerate}
\begin{proof}
\begin{enumerate}[label=(\roman*)]
\setlength{\itemsep}{0pt}
\item Apply Theorem \ref{thm:linear-change-of-variables-nonnegative} to $\abs f$: $\int\abs{f\circ T}\,dm_n=c_T\int\abs f\,dm_n<+\infty$, hence $f\circ T\in\mathscr L^1(m_n)$.
\item Apply Theorem \ref{thm:linear-change-of-variables-nonnegative} to $f^\pm$: $\int f^\pm\circ T\,dm_n=c_T\int f^\pm\,dm_n<+\infty$.\\
Hence $\int f\circ T\,dm_n=\int f^+\circ T\,dm_n-\int f^-\circ T\,dm_n=c_T\int f\,dm_n$.
\end{enumerate}
\hlt{Structural Commentary}: Reduce signed integrals to finite positive parts.\\
\hlt{Key Insight}: Absolute integrability permits subtraction.
\end{proof}
\end{corollary}

\begin{corollary} \label{cor:linear-images-rigid-motion-lebesgue-measure} \hlt{Linear Images and Rigid Motions}\\
Let $E\in\mathcal{L}^n$.
\begin{enumerate}[label=(\roman*)]
\setlength{\itemsep}{0pt}
\item If $T\in\LL(\R^n)$ is singular, then $T(E)$ is Lebesgue measurable and $m_n(T(E))=0$.
\item If $\Psi(x)=Ux+x_0$, where $U$ is a linear isometry, then $\Psi(E)\in\mathcal{L}^n$ and $m_n(\Psi(E))=m_n(E)$.
\end{enumerate}
\begin{proof}
\begin{enumerate}[label=(\roman*)]
\setlength{\itemsep}{0pt}
\item $T(E)\in\mathcal L^n$ by Proposition \ref{prop:lipschitz-images-lebesgue-measurable}. Set $W\defeq\operatorname{range}T$, $H\defeq\R^{n-1}\times\{0\}$; choose invertible $S$ with $S(W)\subseteq H$.\\
Since $H=\bigcup\limits_N(H\cap[-N,N]^n)$ is null, $m_n(W)=0$ by Corollary \ref{cor:invertible-linear-images-lebesgue-sets}; hence $m_n(T(E))=0$.
\item Since $U$ is an isometry, $\abs{\det U}=1$ by Proposition \ref{prop:isometry-determinant-absolute-one}.\\
Thus $U(E)\in\mathcal L^n$ and $m_n(U(E))=m_n(E)$ by Corollary \ref{cor:invertible-linear-images-lebesgue-sets}.\\
By translation invariance, $\Psi(E)=U(E)+x_0\in\mathcal L^n$ and $m_n(\Psi(E))=m_n(U(E))=m_n(E)$.
\end{enumerate}
\hlt{Structural Commentary}: Singular maps land in null subspaces; rigid motions preserve determinant mass.\\
\hlt{Key Insight}: Lebesgue measure records determinant size, not the coordinate description of the set.
\end{proof}
\end{corollary}
